\chapter{Series numéricas}\label{ch:b1:series}

Sumar infinitos números significa tomar el límite de las sumas
parciales --- ni más, ni menos. Este capítulo establece las
definiciones y los criterios de convergencia utilizables en el primer
año: la comparación y los equivalentes para términos positivos, el
\omterm{thm:b1:series:ratio}{criterio del cociente}, la comparación con \omterm{thm:b1:integration:def}{integrales} que da las
\omterm{thm:b1:series:riemann}{series de Riemann}, la \omterm{thm:b1:series:absolute}{convergencia absoluta} y el teorema de las
\omterm{thm:b1:series:alternating}{series alternadas}. La teoría más fina (productos de \omterm{def:b1:series:def}{series},
sumación por paquetes, \omterm{def:b1:series:def}{series} de funciones) corresponde al segundo año.

\section{Generalidades}

\begin{definition}\label{def:b1:series:def}
Dada una sucesión $(u_n)$, la \emph{serie}\index{serie} $\sum u_n$ es
la sucesión de \emph{sumas parciales} $S_N = \sum_{n=0}^{N} u_n$. La
serie \emph{converge} cuando $(S_N)$ converge; el límite es la
\emph{suma} $\sum_{n=0}^{\infty} u_n$, y
$R_N = \sum_{n > N} u_n = S - S_N$ es el \emph{resto}, que tiende a $0$.
\end{definition}

\begin{example}[Series geométricas]\label{ex:b1:series:geometric}
Para $q \in \C$: $\;S_N = \sum_{n=0}^{N} q^n = \frac{1 -
q^{N+1}}{1-q}$ ($q \neq 1$). La \omterm{def:b1:series:def}{serie} converge si y solo si
$\abs q < 1$ (\cref{exo:b1:seq:3}), con\index{serie geométrica}
\[
\sum_{n=0}^{\infty} q^n = \frac{1}{1 - q} .
\]
\end{example}

\begin{example}[Los decimales periódicos son series geométricas]\label{ex:b1:series:periodicdecimal}
¿Qué número es $0.363636\dots$? Su propia escritura es una \omterm{def:b1:series:def}{serie}:
\[
0.\overline{36} = \sum_{k=1}^{\infty} \frac{36}{100^k}
= 36\cdot\frac{1/100}{1 - 1/100} = \frac{36}{99} =
\frac{4}{11} ,
\]
por la suma geométrica con $q = \frac{1}{100}$. En general, un
bloque $B$ de $p$ cifras que se repite para siempre vale
$\frac{B}{10^p - 1}$ --- el mecanismo que hay detrás del criterio de
periodicidad del~\cref{pb:b1:reals:1}, que el lenguaje de este
capítulo enuncia por fin en una línea: un \omterm{pb:b1:reals:1}{desarrollo decimal} es una
\omterm{def:b1:series:def}{serie} convergente, finalmente periódica exactamente cuando su suma
es racional. La maquinaria de cifras del capítulo 10, construida allí
con \omterm{def:b1:reals:bounds}{supremos} desnudos, era teoría de \omterm{def:b1:series:def}{series} viajando de incógnito.
\end{example}

\begin{proposition}[Primeros hechos]\label{prop:b1:series:first}
\begin{enumerate}
  \item Si $\sum u_n$ converge, entonces $u_n \to 0$. (El recíproco
        es \emph{falso}: la \omterm{def:b1:series:def}{serie} armónica.)
  \item Linealidad: las \omterm{def:b1:series:def}{series} convergentes se suman y se
        multiplican por escalares, con las sumas esperadas.
  \item (Telescopaje\index{serie telescópica}) $\sum (v_{n+1} -
        v_n)$ converge si y solo si $(v_n)$ converge, con suma
        $\lim v_n - v_0$.
  \item Cambiar finitos términos no afecta a la convergencia (solo
        a la suma).
\end{enumerate}
\end{proposition}

\begin{proof}
(1) $u_N = S_N - S_{N-1} \to S - S = 0$. La \omterm{def:b1:series:def}{serie} armónica tiene
$u_n = \frac1n \to 0$ y, aun así, diverge (\cref{exo:b1:seq:5}).
(2) Operaciones con límites. (3) $S_N = v_{N+1} - v_0$. (4) Las sumas
parciales cambian en una cantidad finalmente constante.
\end{proof}

\begin{example}[Planificar cifras con el resto geométrico]\label{ex:b1:series:georemainder}
Para $\abs q < 1$, el resto de la \omterm{ex:b1:series:geometric}{serie geométrica} es explícito:
\[
R_N = \sum_{n = N+1}^{\infty} q^n = \frac{q^{N+1}}{1 - q} .
\]
Esto convierte los objetivos de precisión en número de términos antes
de cualquier cálculo. Para evaluar
$\sum_{n\geq0} \bigl(\frac13\bigr)^n = \frac32$ con error $10^{-10}$
hace falta $\frac{(1/3)^{N+1}}{2/3} \leq 10^{-10}$, es decir,
$3^{N} \geq \frac{3}{2}\cdot 10^{10}$, es decir, $N \geq 22$ (pues
$3^{22} \approx 3.1\cdot10^{10}$): veintitrés términos, sabidos de
antemano. Toda estimación de ritmo geométrico de los problemas del
fin de semana (la \omterm{def:b1:series:def}{serie} en $\frac13$ para $\ln 2$, los
\omterm{def:b1:functions:arc}{arcotangentes} de Machin en el~\cref{pb:b1:taylor:1}) es este
presupuesto de dos líneas vestido de gala.
\end{example}

\begin{example}[Un telescopio más largo]\label{ex:b1:series:telescope3}
Calcúlese $\sum_{n\geq1} \frac{1}{n(n+1)(n+2)}$. Fracciones simples
(\cref{ch:b1:fractions}):
\[
\frac{1}{n(n+1)(n+2)}
= \frac{1/2}{n} - \frac{1}{n+1} + \frac{1/2}{n+2}
= \frac12\Bigl(\frac{1}{n(n+1)} - \frac{1}{(n+1)(n+2)}\Bigr),
\]
donde la segunda forma ---una diferencia de valores consecutivos de
$w_n = \frac{1}{n(n+1)}$--- es la telescópica. Por tanto,
\[
\sum_{n=1}^{N} \frac{1}{n(n+1)(n+2)}
= \frac12\Bigl(w_1 - w_{N+1}\Bigr)
= \frac12\Bigl(\frac12 - \frac{1}{(N+1)(N+2)}\Bigr)
\longrightarrow \frac14 .
\]
La idea de cierre: las fracciones simples de tres términos rara vez
telescopan tal como se escriben; reagrúpense antes en una diferencia
$w_n - w_{n+1}$ --- la recompensa no es solo la convergencia, sino la
suma exacta, que ningún criterio de comparación entrega jamás.
\end{example}

\section{Series de términos no negativos}

\begin{theorem}[Sumas parciales acotadas]\label{thm:b1:series:positive}
Si $u_n \geq 0$ para todo $n$, las sumas parciales crecen, de modo
que $\sum u_n$ converge $\iff$ sus sumas parciales están acotadas
superiormente. De ahí el \emph{criterio de comparación}: si
$0 \leq u_n \leq v_n$ para todo $n$ (grande),
\[
\sum v_n \text{ converge} \implies \sum u_n \text{ converge},
\qquad
\sum u_n \text{ diverge} \implies \sum v_n \text{ diverge}.
\]
Y el \emph{criterio de los equivalentes}: si $u_n \sim v_n$ con
$v_n \geq 0$, las dos \omterm{def:b1:series:def}{series} tienen la misma naturaleza.
\end{theorem}

\begin{proof}
Teorema del límite monótono (\cref{thm:b1:seq:monotone}) para lo
primero; comparación de sumas parciales para lo segundo.
Equivalentes: para $n$ grande, $\frac12 v_n \leq u_n \leq 2 v_n$
(definición de $\sim$ con $\varepsilon = \frac12$), y la comparación
se aplica en los dos sentidos.
\end{proof}

\begin{example}[Un equivalente que demuestra la divergencia]\label{ex:b1:series:equivdiverge}
¿Naturaleza de $\sum_{n\geq1} n\sin\dfrac{1}{n^2}$? Como
$\frac{1}{n^2} \to 0$ y $\sin h \sim h$ en $0$:
\[
n\sin\frac{1}{n^2} \;\sim\; n\cdot\frac{1}{n^2} = \frac1n ,
\]
y el criterio de los equivalentes transfiere la divergencia de la
\omterm{def:b1:series:def}{serie} armónica: divergente --- aunque los términos tiendan a $0$.
Un desarrollo, una escala, un veredicto; y el mismo esquema de dos
pasos (equivalente y después consulta de Riemann o geométrica) decide
las cuatro \omterm{def:b1:series:def}{series} del~\cref{exo:b1:series:3}.
\end{example}

\begin{example}[El criterio de los equivalentes en una línea]\label{ex:b1:series:equivtest}
¿Naturaleza de $\sum_{n \geq 1} \frac{\sqrt{n+1} - \sqrt n}{n}$?
Multiplíquese el numerador por el \omterm{def:b1:complex:field}{conjugado}:
\[
\frac{\sqrt{n+1} - \sqrt n}{n}
= \frac{1}{n\,(\sqrt{n+1} + \sqrt n)}
\sim \frac{1}{2\,n^{3/2}} ,
\]
una escala de Riemann convergente ($\alpha = \frac32 > 1$): la
\omterm{def:b1:series:def}{serie} converge. Toda la decisión ha costado un equivalente y una
consulta --- siempre que los términos sean no negativos, como lo son.
La idea de cierre: para las \omterm{def:b1:series:def}{series} positivas, toda la teoría de la
convergencia es un \emph{diccionario de escalas} ($n^{-\alpha}$,
$q^n$, $\frac{1}{n(\ln n)^\alpha}$) más la licencia para sustituir un
término por un equivalente; el trabajo analítico está en la
asintótica (\cref{ch:b1:taylor}), nunca en la sumación.
\end{example}

\begin{theorem}[Comparación con integrales; series de Riemann]\label{thm:b1:series:riemann}
Sea $f$ \omterm{def:b1:continuity:continuous}{continua}, no negativa y \emph{decreciente} en
$\intco{1}{+\infty}$. Entonces
\[
\int_1^{N+1} f(t)\,\dd t \;\leq\; \sum_{n=1}^{N} f(n) \;\leq\; f(1)
+ \int_1^{N} f(t)\,\dd t ,
\]
de modo que $\sum f(n)$ converge si y solo si $\bigl(\int_1^x f\bigr)$
está acotada. En particular, para $\alpha \in \R$:\index{series de Riemann}
\[
\sum_{n \geq 1} \frac{1}{n^\alpha} \text{ converge}
\iff \alpha > 1,
\]
y $\sum_{n=1}^{N} \frac1n = \ln N + O(1)$.
\end{theorem}

\begin{proof}
Para $n \leq t \leq n+1$, la monotonía da
$f(n+1) \leq f(t) \leq f(n)$; integrando en $\intcc{n}{n+1}$ (un
segmento de longitud $1$):
\[
f(n+1) \;\leq\; \int_n^{n+1} f(t)\,\dd t \;\leq\; f(n) .
\]
Sumando las desigualdades de la derecha para $n = 1, \dots, N-1$ se
obtiene $\int_1^{N} f \leq \sum_{n=1}^{N-1} f(n)$, y de ahí el
encuadre superior tras añadir $f(N) \leq f(1)$; sumando las de la
izquierda para $n = 1, \dots, N$ se obtiene
$\sum_{n=2}^{N+1} f(n) \leq \int_1^{N+1} f$, que tras reindexar es el
encuadre inferior. Convergencia: las sumas parciales y las
\omterm{thm:b1:integration:def}{integrales} $\int_1^x f$ se acotan mutuamente salvo la constante
$f(1)$, y las dos son no decrecientes, de modo que una está acotada
si y solo si lo está la otra (\cref{thm:b1:series:positive}). Para
$f(t) = t^{-\alpha}$ ($\alpha \neq 1$):
$\int_1^x t^{-\alpha}\dd t = \frac{x^{1-\alpha} - 1}{1 - \alpha}$,
acotada si y solo si $\alpha > 1$; para $\alpha = 1$, la \omterm{thm:b1:integration:def}{integral} es
$\ln x \to \infty$, y el encuadre da $\ln(N+1) \leq H_N \leq 1 +
\ln N$. Y para $\alpha \leq 0$ los términos no tienden a $0$.
\end{proof}

\begin{example}[La pila armónica]\label{ex:b1:series:harmonicstack}
¿Cuántos términos tiene que acumular la \omterm{def:b1:series:def}{serie} armónica para pasar de
$20$? El encuadre $\ln(N+1) \leq H_N \leq 1 + \ln N$ responde sin
sumar nada: $H_N \geq 20$ exige $1 + \ln N \geq 20$, es decir,
$N \geq \eu^{19} \approx 1.8\cdot10^{8}$, y está garantizado en
cuanto $\ln(N + 1) \geq 20$, es decir, $N \approx \eu^{20} \approx
4.9\cdot10^{8}$. (El problema del fin de semana lo afina hasta
$N \approx \eu^{20 - \gamma} \approx 2.7\cdot10^{8}$ mediante la
\omterm{pb:b1:series:1}{constante de Euler}.)
La idea de cierre: la comparación con \omterm{thm:b1:integration:def}{integrales} no solo decide la
convergencia --- también \emph{localiza} las sumas parciales con
precisión logarítmica, convirtiendo un cálculo desesperado (cientos
de millones de términos) en una estimación de dos líneas.
\end{example}

\begin{theorem}[Criterio del cociente (d'Alembert)]\label{thm:b1:series:ratio}
Sea $u_n > 0$ con $\frac{u_{n+1}}{u_n} \to \ell$.\index{criterio del cociente}
\begin{itemize}
  \item Si $\ell < 1$: $\sum u_n$ converge;
  \item si $\ell > 1$: $u_n \to +\infty$, divergencia;
  \item si $\ell = 1$: no hay conclusión ($\sum \frac1n$ diverge y
        $\sum \frac{1}{n^2}$ converge).
\end{itemize}
\end{theorem}

\begin{proof}
Si $\ell < 1$, fíjese $q \in \intoo{\ell}{1}$: más allá de cierto $N$,
$u_{n+1} \leq q\,u_n$, luego $u_n \leq u_N q^{\,n-N}$ por inducción:
comparación con una \omterm{ex:b1:series:geometric}{serie geométrica}. Si $\ell > 1$: más allá de
cierto $N$ la sucesión $(u_n)$ es creciente, de modo que no puede
tender a $0$ (su límite, si lo tiene, es $\geq u_N > 0$); por
la~\cref{prop:b1:series:first}~(1), divergencia --- y, de hecho,
$u_n \geq u_N q^{n-N}$ con $q > 1$ da $u_n \to \infty$.
\end{proof}

\begin{example}\label{ex:b1:series:ratio}
$\sum \frac{x^n}{n!}$ converge para todo $x > 0$: cociente
$\frac{x}{n+1} \to 0$. Su suma es $\eu^x$: por Taylor--Lagrange
(\cref{thm:b1:taylor:lagrange}) en $\intcc{0}{x}$,
\[
\Bigl| \eu^x - \sum_{k=0}^{n} \frac{x^k}{k!} \Bigr|
\leq \eu^{x}\, \frac{x^{n+1}}{(n+1)!} \xrightarrow[n\to\infty]{} 0 ,
\]
tendiendo la cota a $0$ porque el factorial domina
(\cref{exo:b1:integration:9}~(1) usó el mismo hecho). El mismo
argumento suma las \omterm{def:b1:series:def}{series} de $\sin$, $\cos$, $\sinh$ y $\cosh$ en
todo $\R$.
\end{example}

\begin{example}[El criterio del cociente es suficiente, no necesario]\label{ex:b1:series:rationotnecessary}
Sea $u_n = 2^{-n}$ para $n$ par y $u_n = 2^{-n-2}$ para $n$ impar.
Los cocientes consecutivos oscilan entre $\frac{1}{8}$ y
$\frac12\cdot4 = 2$, de modo que $\frac{u_{n+1}}{u_n}$ no tiene
límite y d'Alembert enmudece --- y, sin embargo, $u_n \leq 2^{-n}$ y
el criterio de comparación resuelve la convergencia al instante. La
hipótesis del criterio (que el cociente \emph{converja}) es una
restricción real: le van bien los términos con una única estructura
multiplicativa dominante (factoriales, potencias) y falla con
cualquier cosa que respire. Cuando los cocientes se portan mal,
retrocédase a la comparación con una envolvente geométrica --- que es
todo lo que el \omterm{thm:b1:series:ratio}{criterio del cociente} fue siempre, como muestra su
demostración.
\end{example}

\begin{example}[El criterio del cociente en batallas de factoriales]\label{ex:b1:series:ratiobinom}
¿Naturaleza de $\sum_{n\geq0} \dfrac{(n!)^2}{(2n)!}$ (los inversos de
los \omterm{def:b1:counting:objects}{coeficientes binomiales} centrales, salvo el factor $n + 1$)? El
cociente hace colapsar los factoriales:
\[
\frac{u_{n+1}}{u_n}
= \frac{((n+1)!)^2}{(n!)^2}\cdot\frac{(2n)!}{(2n+2)!}
= \frac{(n+1)^2}{(2n+1)(2n+2)}
\longrightarrow \frac14 < 1 :
\]
convergente, y con holgura --- los términos decrecen esencialmente
como $4^{-n}$, de acuerdo con
$\binom{2n}{n} \geq \frac{4^n}{2n+1}$ del~\cref{pb:b1:integration:1}.
La idea de cierre: los cocientes de factoriales son exactamente lo
que el \omterm{thm:b1:series:ratio}{criterio del cociente} digiere --- todo factorial se cancela
dejando una función racional de $n$, cuyo límite se lee en los
términos dominantes.
\end{example}

\section{Convergencia absoluta; series alternadas}

\begin{theorem}[Convergencia absoluta]\label{thm:b1:series:absolute}
Si $\sum \abs{u_n}$ converge (\emph{convergencia
absoluta}\index{convergencia absoluta}), entonces $\sum u_n$
converge, y $\bigl|\sum u_n\bigr| \leq \sum \abs{u_n}$. Esto vale
para términos reales o complejos.
\end{theorem}

\begin{proof}
Las sumas parciales cumplen, para $M > N$ (criterio de Cauchy,
\cref{thm:b1:seq:complete}):
\[
\abs{S_M - S_N} = \Bigl| \sum_{n=N+1}^{M} u_n \Bigr|
\leq \sum_{n=N+1}^{M} \abs{u_n},
\]
que es pequeño para $N$ grande, ya que las sumas parciales de
$\sum\abs{u_n}$ forman una \omterm{def:b1:seq:cauchy}{sucesión de Cauchy}. Luego $(S_N)$ es de
Cauchy y, por tanto, converge. La desigualdad pasa al límite desde la
desigualdad triangular finita.
\end{proof}

\begin{example}[Convergencia absoluta, real y compleja]\label{ex:b1:series:absexamples}
$\sum_{n\geq1} \frac{\sin n}{n^2}$: los términos cambian de signo de
forma errática (en efecto, $(\sin n)$ es \omterm{def:b1:topology:dense}{denso} en $\intcc{-1}{1}$,
\cref{exo:b1:seq:12}) y no se ve ninguna estructura alternada. La
\omterm{thm:b1:series:absolute}{convergencia absoluta} lo rescata todo de golpe:
$\bigl|\frac{\sin n}{n^2}\bigr| \leq \frac{1}{n^2}$, una escala
convergente, luego la \omterm{def:b1:series:def}{serie} converge. El mismo escudo funciona
sobre $\C$: $\sum_{n\geq1}\frac{\eu^{\iu n}}{n^2}$ converge porque
$\bigl|\frac{\eu^{\iu n}}{n^2}\bigr| = \frac{1}{n^2}$ --- los
patrones de signos, incluso los bidimensionales, son irrelevantes en
cuanto los \omterm{def:b1:complex:field}{módulos} son sumables. La idea de cierre: la
\omterm{thm:b1:series:absolute}{convergencia absoluta} es la única herramienta de este capítulo que
nunca pregunta cómo están organizados los signos; pruébese primero
(\cref{met:b1:series:decide}) y resérvense los criterios delicados
para las \omterm{def:b1:series:def}{series} que no la cumplen.
\end{example}

\begin{theorem}[Criterio de las series alternadas]\label{thm:b1:series:alternating}
Sea $(a_n)$ decreciente con $a_n \to 0$. Entonces la \omterm{def:b1:series:def}{serie} alternada
$\sum (-1)^n a_n$ converge; su suma está entre dos sumas parciales
consecutivas cualesquiera, y\index{series alternadas}
\[
\abs{R_N} = \Bigl| \sum_{n > N} (-1)^n a_n \Bigr| \leq a_{N+1} .
\]
\end{theorem}

\begin{proof}
Las sumas parciales pares e impares son adyacentes:
$S_{2p+2} - S_{2p} = a_{2p+2} - a_{2p+1} \leq 0$ (decrecimiento),
$S_{2p+1} - S_{2p-1} = a_{2p} - a_{2p+1} \geq 0$ (crecimiento) y
$S_{2p} - S_{2p+1} = a_{2p+1} \to 0$. Por
el~\cref{thm:b1:seq:adjacent} comparten un límite $S$, que el
criterio de las dos subsucesiones
(\cref{prop:b1:seq:subsequences}) convierte en el límite de $(S_N)$;
además, $S$ queda atrapado entre sumas parciales consecutivas, y
$\abs{S - S_N}$ es a lo sumo la distancia a la siguiente, $a_{N+1}$.
\end{proof}

\begin{example}[Serie armónica alternada]\label{ex:b1:series:ln2}
$\sum_{n \geq 1} \frac{(-1)^{n-1}}{n}$ converge (criterio de las
alternadas) pero no absolutamente (\omterm{def:b1:series:def}{serie} armónica). Su suma es
$\ln 2$: a partir de la identidad geométrica finita
$\frac{1}{1+t} = \sum_{k=0}^{n-1} (-t)^k + \frac{(-t)^n}{1+t}$,
intégrese en $\intcc{0}{1}$:
\[
\ln 2 = \sum_{k=1}^{n} \frac{(-1)^{k-1}}{k}
+ (-1)^n \int_0^1 \frac{t^n}{1+t}\,\dd t ,
\qquad
0 \leq \int_0^1 \frac{t^n}{1+t}\,\dd t \leq \frac{1}{n+1} \to 0 .
\]
La convergencia es dolorosamente lenta ($R_N \approx \frac{1}{N}$) ---
las \omterm{thm:b1:series:alternating}{series alternadas} convergen por cancelación, no por pequeñez.
\end{example}

\begin{omfigure}
\begin{tikzpicture}[x=0.78cm, y=7.2cm]
  \draw[-Stealth] (0.3,0.42) -- (11,0.42)
    node[below, font=\small] {$N$};
  \draw (1,0.415) -- (1,0.425) node[below=1pt, font=\small]
    {$1$};
  \draw (5,0.415) -- (5,0.425) node[below=1pt, font=\small]
    {$5$};
  \draw (10,0.415) -- (10,0.425) node[below=1pt, font=\small]
    {$10$};
  \draw[gray, dashed] (0.3,0.6931) -- (10.7,0.6931)
    node[right, font=\small, text=black] {$\ln 2$};
  \node[left, font=\small] at (0.3,1) {$1$};
  \node[left, font=\small] at (0.3,0.5) {$0.5$};
  \draw[omDef, thick]
    plot coordinates {(1,1) (2,0.5) (3,0.8333) (4,0.5833)
    (5,0.7833) (6,0.6167) (7,0.7595) (8,0.6345)
    (9,0.7456) (10,0.6456)};
  \fill[omDef] (1,1) circle (1.5pt);
  \fill[omDef] (2,0.5) circle (1.5pt);
  \fill[omDef] (3,0.8333) circle (1.5pt);
  \fill[omDef] (4,0.5833) circle (1.5pt);
  \fill[omDef] (5,0.7833) circle (1.5pt);
  \fill[omDef] (6,0.6167) circle (1.5pt);
  \fill[omDef] (7,0.7595) circle (1.5pt);
  \fill[omDef] (8,0.6345) circle (1.5pt);
  \fill[omDef] (9,0.7456) circle (1.5pt);
  \fill[omDef] (10,0.6456) circle (1.5pt);
\end{tikzpicture}

{\small Las sumas parciales $S_N$ de la \omterm{def:b1:series:def}{serie} armónica alternada
$1 - \frac12 + \frac13 - \cdots$ saltan por encima de su límite
$\ln 2$ en cada paso: las sumas impares desde arriba y las pares
desde abajo, con saltos de tamaño $\frac{1}{N+1}$. El encuadre es la
demostración del~\cref{thm:b1:series:alternating} hecha visible --- y
el lento cierre de la tenaza ($\abs{S_N - \ln 2} \approx
\frac{1}{2N}$, problema del fin de semana \cref{pb:b1:series:1}) es
la razón de que nadie calcule $\ln 2$ así.}
\end{omfigure}

\begin{remark}[Errores frecuentes con las series]
(i) \emph{El criterio de los equivalentes necesita un signo}: sean
$v_n = \frac{(-1)^n}{\sqrt n}$ y $u_n = v_n + \frac1n$. Entonces
$\frac{u_n}{v_n} = 1 + \frac{(-1)^n}{\sqrt n} \to 1$, luego
$u_n \sim v_n$; y, sin embargo, $\sum v_n$ converge (criterio de las
alternadas) mientras que $\sum u_n = \sum v_n + \sum \frac1n$
diverge. La equivalencia controla el \emph{tamaño} de los términos y,
para las \omterm{def:b1:series:def}{series} con signos, el tamaño no es el destino --- el
criterio está \omterm{def:b1:logic:statement}{enunciado}, y es cierto, solo para términos (finalmente)
no negativos. (ii) \emph{$u_n \to 0$ no demuestra nada}: la \omterm{def:b1:series:def}{serie}
armónica es el contraejemplo eterno; la dirección recíproca
(\cref{prop:b1:series:first}~(1)) es solo un criterio rápido de
divergencia. (iii) \emph{Un cociente con límite $1$ es silencio, no
convergencia}: tanto $\sum\frac1n$ como $\sum\frac{1}{n^2}$ tienen
cociente $\to 1$; pásese a las escalas de Riemann o a la comparación
con \omterm{thm:b1:integration:def}{integrales}. (iv) \emph{Lo alternado necesita ser decreciente}:
$\sum \frac{(-1)^n}{n + (-1)^n}$ parece alternada y solo se trata con
un desarrollo (\cref{exo:b1:series:5}); el problema del fin de semana
del~\cref{ch:b1:taylor} (su pregunta 23) muestra que el criterio
puede fallar de plano sin monotonía. (v) \emph{Agrupar y reordenar no
salen gratis}: poner paréntesis es inofensivo para las \omterm{def:b1:series:def}{series}
convergentes, pero puede crear convergencia a partir de divergencia
($1 - 1 + 1 - \cdots$ agrupado por parejas), y reordenar puede
cambiar la propia suma --- el drama que escenifica el problema del
fin de semana de este capítulo (\cref{pb:b1:series:1}).
\end{remark}

\begin{method}[Decidir la naturaleza de una serie]\label{met:b1:series:decide}
\begin{enumerate}
  \item ¿Es $u_n \to 0$? Si no, divergencia; alto.
  \item Términos no negativos: búsquese un equivalente de $u_n$
        (¡desarrollos, \cref{ch:b1:taylor}!) y compárese con las
        escalas de Riemann o geométricas; los factoriales y las
        potencias piden el \omterm{thm:b1:series:ratio}{criterio del cociente}; y una $f(n)$
        decreciente pide la comparación con \omterm{thm:b1:integration:def}{integrales}.
  \item Los signos varían: pruébese primero la \omterm{thm:b1:series:absolute}{convergencia
        absoluta}; si falla, el criterio de las alternadas
        (compruébese con cuidado el \emph{decrecimiento}); y más
        allá, herramientas de segundo año.
\end{enumerate}
\end{method}

\begin{example}[Denominadores impares, medio telescopio]\label{ex:b1:series:oddtelescope}
Calcúlese $\sum_{n \geq 1} \dfrac{1}{4n^2 - 1}$. Fracciones simples:
$\frac{1}{(2n-1)(2n+1)} = \frac12\bigl(\frac{1}{2n-1} -
\frac{1}{2n+1}\bigr)$, luego
\[
\sum_{n=1}^{N} \frac{1}{4n^2 - 1}
= \frac12\Bigl(1 - \frac{1}{2N+1}\Bigr)
\longrightarrow \frac12 .
\]
Compárese con $\sum \frac{1}{n(n+1)} = 1$
(\cref{exo:b1:series:1}): el mismo esqueleto telescópico, pero aquí
los términos consecutivos están a distancia dos en los impares, y el
factor $\frac12$ registra el paso. La idea de cierre: telescopar es
un cambio de punto de vista, no un truco --- siempre que el término
general sea una diferencia $w_n - w_{n+1}$ de una sucesión con
límite, la suma es $w_1 - \lim w$, exactamente
la~\cref{prop:b1:series:first}~(3).
\end{example}

\begin{remark}[La cadena del análisis, en retrospectiva]
Este capítulo es donde converge el análisis del volumen, y cada
criterio nombra a su antepasado. Las sumas parciales acotadas son el
teorema del límite monótono (\cref{ch:b1:seq}), que a su vez es el
axioma de completitud del~\cref{ch:b1:reals}; la \omterm{thm:b1:series:absolute}{convergencia
absoluta} es el criterio de Cauchy; el criterio de la \omterm{thm:b1:integration:def}{integral} es el
encuadre de áreas del~\cref{ch:b1:integration}; los equivalentes de
los términos generales son los desarrollos del~\cref{ch:b1:taylor}; y
el teorema de las alternadas es el lema de las \omterm{thm:b1:seq:adjacent}{sucesiones adyacentes}
con su traje de domingo. Leída al revés, la cadena explica para
\emph{qué} servía cada capítulo --- y los problemas del fin de semana
que la enhebran (las cifras $b$-ádicas, Cesàro--Stolz, las máquinas de
irracionalidad, la \omterm{pb:b1:series:1}{constante de Euler}) son las mismas pocas ideas
encontrándose a altitudes cada vez mayores. El álgebra lineal que
sigue cambia de asunto, no de nivel de exigencia: la costumbre de los
\omterm{def:b1:logic:statement}{enunciados} exactos con error certificado sobrevive al paso de los
límites a las dimensiones.
\end{remark}

\begin{remark}[Adónde van las series a continuación]
Este capítulo cierra el análisis del volumen y abre tres puertas. En
el volumen del segundo año, las \omterm{def:b1:series:def}{series} adquieren una variable
($\sum a_n x^n$: \omterm{def:b1:series:def}{series} de potencias, con su radio de convergencia) y
después una teoría con valores funcionales (\omterm{def:b1:series:def}{series} de Fourier); y la
dicotomía \omterm{thm:b1:series:absolute}{convergencia absoluta} frente a condicional, dramatizada en
el problema del fin de semana de más abajo, se convierte en la piedra
angular de las dos. En probabilidad (volumen del tercer año), las
esperanzas de las variables aleatorias discretas \emph{son} \omterm{def:b1:series:def}{series}, y
la \omterm{thm:b1:series:absolute}{convergencia absoluta} es lo que las hace estar bien definidas. Y
la \omterm{def:b1:series:def}{serie} de Riemann $\sum n^{-s}$, llevada a $s$ complejo, se
convierte en la función zeta --- la \omterm{def:b1:series:def}{serie} más estudiada de las
matemáticas.
\end{remark}

\section{Ejercicios}

\begin{exercise}[$\star$]\label{exo:b1:series:1}
Naturaleza (y suma, cuando telescope) de:
\[
\sum_{n\geq1} \frac{1}{n(n+1)},
\qquad
\sum_{n\geq2} \ln\Bigl(1 - \frac{1}{n^2}\Bigr),
\qquad
\sum_{n\geq0} \frac{3^n + 4^n}{5^n} .
\]
\end{exercise}

\begin{exercise}[$\star$]\label{exo:b1:series:2}
Naturaleza de: $\;\sum \dfrac{n^2}{2^n}$; $\;\sum \dfrac{n!}{n^n}$;
$\;\sum \dfrac{2^n\,n!}{n^n}$; $\;\sum \dfrac{3^n\,n!}{n^n}$.
\emph{(\omterm{thm:b1:series:ratio}{Criterio del cociente}; recuérdese $\bigl(1 + \frac1n\bigr)^n \to \eu$.)}
\end{exercise}

\begin{exercise}[$\star$]\label{exo:b1:series:3}
Naturaleza de: $\;\sum \sin\dfrac{1}{n^2}$; $\;\sum
\Bigl(1 - \cos\dfrac1n\Bigr)$; $\;\sum \dfrac{1}{\sqrt{n(n+1)}}$;
$\;\sum \dfrac{\ln n}{n^2}$ \emph{(compárese con $n^{-3/2}$)}.
\end{exercise}

\begin{exercise}[$\star$]\label{exo:b1:series:4}
Demuéstrese que $\sum_{n\geq1} \frac{1}{n^2}$ converge con suma
$\leq 2$, usando $\frac{1}{n^2} \leq \frac{1}{n(n-1)}$ para
$n \geq 2$ y una cota telescópica.
\end{exercise}

\begin{exercise}[$\star\star$]\label{exo:b1:series:5}
Naturaleza de $\;\sum \dfrac{(-1)^n}{\sqrt n}$, de $\;\sum
\dfrac{(-1)^n}{n + (-1)^n}$ \emph{(desarróllese: el criterio de las
alternadas no se aplica directamente --- ¿por qué?)} y de $\;\sum
\sin\bigl(\pi\sqrt{n^2+1}\,\bigr)$ \emph{(redúzcase \omterm{def:b1:complex:field}{módulo} $\pi$:
$\sqrt{n^2+1} = n + \frac{1}{2n} + O(n^{-3})$)}.
\end{exercise}

\begin{exercise}[$\star\star$]\label{exo:b1:series:6}
¿Para qué $\alpha > 0$ converge
$\sum_{n \geq 2} \dfrac{1}{n (\ln n)^{\alpha}}$? \emph{(Comparación
con \omterm{thm:b1:integration:def}{integrales}; sustitúyase $u = \ln t$.)}
\end{exercise}

\begin{exercise}[$\star\star$]\label{exo:b1:series:7}
Sea $u_n = \dfrac{1}{n} - \ln\Bigl(1 + \dfrac1n\Bigr)$. Demuéstrese
que $0 \leq u_n \leq \dfrac{1}{2n^2}$, que $\sum u_n$ converge y
dedúzcase la existencia de la \omterm{pb:b1:series:1}{constante de Euler}\index{constante de Euler}:
\[
\gamma = \lim_{N \to \infty} \Bigl( \sum_{n=1}^{N} \frac 1n - \ln N
\Bigr) .
\]
\end{exercise}

\begin{exercise}[$\star\star$]\label{exo:b1:series:8}
Calcúlense las sumas
\[
\sum_{n=1}^{\infty} \frac{1}{n(n+2)}
\qquad\text{y}\qquad
\sum_{n=0}^{\infty} \frac{n}{2^n} .
\]
\emph{(Para la primera: fracciones simples. Para la segunda:
calcúlese $\sum_{n=1}^{N} n x^{n-1}$ en forma \omterm{def:b1:topology:closed}{cerrada} y hágase
$N \to \infty$ en $x = \frac12$.)}
\end{exercise}

\begin{exercise}[$\star\star\star$]\label{exo:b1:series:9}
(Condensación de Cauchy) Sea $(u_n)$ no negativa y decreciente.
Demuéstrese que
\[
\sum_{n \geq 1} u_n \text{ converge}
\iff
\sum_{k \geq 0} 2^k\, u_{2^k} \text{ converge},
\]
comparando paquetes de términos entre potencias consecutivas de $2$.
Recupérense a partir de ahí el criterio de Riemann y
el~\cref{exo:b1:series:6}.
\end{exercise}

\begin{exercise}[$\star\star\star$]\label{exo:b1:series:10}
Usando la identidad \omterm{thm:b1:integration:def}{integral} del~\cref{ex:b1:series:ln2} adaptada a
$\frac{1}{1+t^2}$, demuéstrese la fórmula de Leibniz
\[
\frac{\pi}{4} = \sum_{n=0}^{\infty} \frac{(-1)^n}{2n+1}
= 1 - \frac13 + \frac15 - \frac17 + \cdots
\]
con la cota de error $\abs{R_N} \leq \frac{1}{2N+3}$.
\end{exercise}

\begin{exercise}[$\star\star$]\label{exo:b1:series:11}
Naturaleza de $\displaystyle\sum_{n \geq 1} \frac{1}{n^{1 + 1/n}}$.
\emph{(Calcúlese el límite de $n^{1/n}$ y hállese un equivalente del
término general: el criterio de Riemann necesita un exponente
\emph{fijo}.)}
\end{exercise}

\begin{exercise}[$\star\star\star$]\label{exo:b1:series:12}
Sea $(u_n)$ no negativa y \emph{decreciente} con $\sum u_n$
convergente. Demuéstrese que $n\,u_n \to 0$ \emph{(acótese
$n\,u_{2n}$ por una rodaja $\sum_{k=n+1}^{2n} u_k$ y úsese el
criterio de Cauchy)}. Véase que el recíproco falla y que la hipótesis
de monotonía no se puede suprimir.
\end{exercise}

\section{Problema: la constante de Euler y la serie que cambia de
suma}

\begin{problem}[{Problema del fin de semana --- $H_n = \ln n + \gamma +
\frac{1}{2n} + O(n^{-2})$, y la reordenación de
$1 - \frac12 + \frac13 - \dots$ hasta $\frac{\ln 2}{2}$}]
\label{pb:b1:series:1}
Dos historias comparten la \omterm{def:b1:series:def}{serie} armónica. Primero, la contabilidad
exacta de su divergencia: $H_n - \ln n$ converge a la \omterm{pb:b1:series:1}{constante de
Euler} $\gamma$ (\cref{exo:b1:series:7}), y este problema afina el
\omterm{def:b1:logic:statement}{enunciado} hasta una ley con dos lados
$\frac{1}{2(n+1)} \leq H_n - \ln n - \gamma \leq \frac{1}{2n}$, que
certifica $\gamma = 0.5772\dots$ a mano. Segundo, el escándalo de la
convergencia condicional: la \omterm{def:b1:series:def}{serie} armónica alternada suma $\ln 2$
(\cref{ex:b1:series:ln2}) y, sin embargo, \emph{los mismos términos,
en otro orden}, suman $\frac{\ln 2}{2}$ ---o
$\ln 2 + \frac12\ln\frac pq$ para cualesquiera $p, q$, o cualquier
real que se quiera (Riemann)---. Las dos historias son una sola: las
sumas reordenadas se calculan \emph{con} la ley de
$\gamma$.\index{constante de Euler}\index{reordenación de series}

\medskip
\noindent\textbf{Parte I --- $\gamma$, encuadrada.} Póngase
$a_n = H_n - \ln n$ y $b_n = H_n - \ln(n+1)$.
\begin{enumerate}
  \item Usando $\frac{t}{1+t} \leq \ln(1 + t) \leq t$, véase que
        $(a_n)$ decrece, que $(b_n)$ crece y que son adyacentes; su
        límite común es $\gamma$, con $b_n \leq \gamma \leq a_n$
        para todo $n$.
  \item Primer disparo numérico: a partir de
        $H_{10} = 2.928968\dots$, encuádrese $\gamma$ entre
        $b_{10} = 0.5311$ y $a_{10} = 0.6264$. ¿Qué $n$ haría falta
        para cuatro decimales con este encuadre tosco?
  \item Véase la representación exacta de la cola
        $a_n - \gamma = \sum_{k \geq n} w_k$ (límite de sumas
        parciales), donde
        \[
        w_k = a_k - a_{k+1}
        = \ln\Bigl(1 + \frac1k\Bigr) - \frac{1}{k+1}
        = \int_k^{k+1} \frac{(k + 1 - t)}{t\,(k+1)}\,\dd t ,
        \]
        y dedúzcase de la forma integral la cota por los dos lados
        $\dfrac{1}{2(k+1)^2} \leq w_k \leq \dfrac{1}{2k(k+1)}$.
\end{enumerate}

\medskip
\noindent\textbf{Parte II --- La ley del $\frac{1}{2n}$.}
\begin{enumerate}[resume]
  \item Súmense las cotas de la pregunta 3 (los dos lados telescopan
        o se comparan con telescopios) y conclúyase la ley:
        \[
        \frac{1}{2(n+1)} \;\leq\; H_n - \ln n - \gamma \;\leq\;
        \frac{1}{2n} \qquad (n \geq 1).
        \]
  \item Dedúzcase $H_n = \ln n + \gamma + \frac{1}{2n} +
        O\bigl(\frac{1}{n^2}\bigr)$; con precisión, véase que
        $\gamma_n = H_n - \ln n - \frac{1}{2n}$ cumple
        $-\frac{1}{2n(n+1)} \leq \gamma_n - \gamma \leq 0$.
  \item Certifíquense cuatro decimales con $n = 100$: dado
        $H_{100} = 5.1873775\dots$, calcúlese
        $\gamma_{100} = 0.577207\dots$ y conclúyase
        $\gamma = 0.5772 \pm 5\cdot10^{-5}$ (valor verdadero
        $0.5772156\dots$).
  \item Dos dividendos de la ley, los dos necesarios más adelante:
        cuando $m \to \infty$,
        \[
        H_{2m} - H_m = \ln 2 - \frac{1}{4m} +
        O\Bigl(\frac{1}{m^2}\Bigr),
        \qquad
        \sum_{j=1}^{m} \frac{1}{2j-1} = \frac{\ln m}{2} + \ln 2
        + \frac\gamma2 + o(1) ,
        \]
        el segundo mediante
        $\sum_{j \leq m} \frac{1}{2j-1} = H_{2m} - \frac12 H_m$, y
        análogamente
        $\sum_{j=1}^{m} \frac{1}{2j} = \frac{\ln m}{2} +
        \frac\gamma2 + o(1)$.
\end{enumerate}

\medskip
\noindent\textbf{Parte III --- La \omterm{def:b1:series:def}{serie} armónica alternada, hasta el
segundo orden.}
\begin{enumerate}[resume]
  \item Véase (por inducción, o agrupando) la identidad
        $\sum_{k=1}^{2m} \frac{(-1)^{k-1}}{k} = H_{2m} - H_m$, y
        dedúzcanse tanto la suma $\ln 2$ (otra vez) como la
        velocidad exacta:
        \[
        \sum_{k=1}^{2m} \frac{(-1)^{k-1}}{k}
        = \ln 2 - \frac{1}{4m} + O\Bigl(\frac{1}{m^2}\Bigr) .
        \]
  \item Dedúzcase el error asintótico de la \omterm{def:b1:series:def}{serie} armónica alternada
        en \emph{cualquier} índice: $S - S_N \sim
        \frac{(-1)^N}{2N}$ --- dos veces menor que la cota del peor
        caso $a_{N+1} \approx \frac1N$
        del~\cref{thm:b1:series:alternating}.
  \item (Aceleración gratis) Véase que las sumas promediadas
        $\tilde S_N = \frac{S_N + S_{N+1}}{2}$ cumplen
        $\tilde S_N = \ln 2 + O\bigl(\frac{1}{N^2}\bigr)$.
        Compruébese: $S_{10} = 0.64563$, $S_{11} = 0.73654$,
        $\tilde S_{10} = 0.69109$, frente a $\ln 2 = 0.69315$: un
        promedio compra dos decimales.
  \item Explíquese en dos frases por qué ningún truco así puede
        ayudar a un fenómeno de cola \emph{positiva} y divergente
        como el encuadre de la pregunta 2: el error alternado
        oscila (signo $(-1)^N$), de modo que promediar cancela su
        término dominante, mientras que el error
        $\frac{1}{2n}$ del encuadre de $\gamma$ tiene signo
        constante. (Promediar $a_n$ y $b_n$ \emph{sí} ayuda:
        relaciónese $\frac{a_n + b_n}{2}$ con la estimación del
        punto medio $H_n - \ln\bigl(n + \frac12\bigr)$ y véase que
        su error es $O\bigl(\frac{1}{n^2}\bigr)$.)
\end{enumerate}

\medskip
\noindent\textbf{Parte IV --- La rigidez y su fracaso.}
\begin{enumerate}[resume]
  \item Véase que la parte positiva $\sum \frac{1}{2j-1}$ y la parte
        negativa $\sum \frac{1}{2j}$ de la \omterm{def:b1:series:def}{serie} armónica alternada
        divergen las dos --- la firma de la convergencia
        \emph{condicional}.
  \item Demuéstrese el \omterm{def:b1:logic:statement}{enunciado} general que hay detrás de la
        pregunta 12: si $\sum u_n$ converge pero $\sum \abs{u_n}$
        diverge, entonces la \omterm{def:b1:series:def}{serie} de las partes positivas
        $\sum u_n^+$ y la de las partes negativas $\sum u_n^-$
        divergen las dos \emph{(a partir de
        $u_n^\pm = \frac{\abs{u_n} \pm u_n}{2}$: si una convergiera,
        también la otra, y por tanto $\sum\abs{u_n}$)}. Ese
        reservorio inagotable de masa positiva y negativa es lo que
        gastará la receta de Riemann.
  \item (Rigidez) Demuéstrese: si $\sum u_n$ converge
        \emph{absolutamente} y $\sigma \colon \N \to \N$ es una
        \omterm{def:b1:logic:inj}{biyección}, entonces $\sum u_{\sigma(n)}$ converge a la misma
        suma \emph{(para $N$ grande, los primeros $M$ términos
        reordenados contienen $u_0, \dots, u_N$; compárense las
        sumas parciales a través de la cola
        $\sum_{n > N}\abs{u_n}$)}.
  \item (La receta de Riemann) Sea $t \in \R$. Descríbase la
        reordenación voraz de la \omterm{def:b1:series:def}{serie} armónica alternada: tómense
        términos positivos $1, \frac13, \frac15, \dots$ hasta que la
        suma parcial supere por primera vez $t$, después términos
        negativos $-\frac12, -\frac14, \dots$ hasta que baje por
        primera vez de $t$, y repítase. Véase que cada término se
        usa exactamente una vez, que tras el primer cruce las sumas
        parciales se quedan a distancia de $t$ menor que el último
        término usado, y conclúyase que la \omterm{def:b1:series:def}{serie} reordenada
        converge a $t$: \emph{cualquier} suma prescrita es alcanzable.
\end{enumerate}

\medskip
\noindent\textbf{Parte V --- La fórmula $(p, q)$.} Fíjense enteros
$p, q \geq 1$. Reordénese la \omterm{def:b1:series:def}{serie} armónica alternada en bloques:
$p$ términos positivos (los siguientes inversos impares), después $q$
términos negativos (los siguientes inversos pares), y repítase.
\begin{enumerate}[resume]
  \item Compruébese que se trata de una reordenación genuina (cada
        término exactamente una vez) y que para $(p, q) = (1, 2)$ dice
        \[
        1 - \frac12 - \frac14 + \frac13 - \frac16 - \frac18 +
        \frac15 - \frac1{10} - \frac1{12} + \dots
        \]
  \item (La bisección exacta) Para $(p, q) = (1, 2)$, demuéstrese la
        identidad de bloque
        \[
        \frac{1}{2k-1} - \frac{1}{4k-2} - \frac{1}{4k}
        = \frac12\Bigl(\frac{1}{2k-1} - \frac{1}{2k}\Bigr),
        \]
        y dedúzcase la relación \emph{exacta}
        $T_{3K} = \frac12 S_{2K}$ entre las sumas parciales
        reordenadas y las originales: la reducción a la mitad de la
        suma es visible en cada etapa finita, no solo en el límite.
  \item Véase que la suma parcial tras $K$ bloques completos vale
        $\sum_{j=1}^{pK} \frac{1}{2j-1} - \sum_{j=1}^{qK}
        \frac{1}{2j}$, y calcúlese su límite con la pregunta 7:
        \[
        \ln 2 + \frac12 \ln\frac pq .
        \]
  \item Contrólense las sumas parciales \emph{dentro} de un bloque
        (los términos tienden a $0$) y conclúyase que la \omterm{def:b1:series:def}{serie}
        reordenada según $(p, q)$ converge a
        $\ln 2 + \frac12\ln \frac pq$. En particular, $(1, 2)$ da
        $\frac{\ln 2}{2}$: compruébese con los nueve primeros
        términos, $T_9 = 0.3083$, avanzando hacia $0.3466$.
  \item Comprobaciones de coherencia y alcance: $(1,1)$ recupera
        $\ln 2$; $(2,1)$ da $\frac32\ln 2$; ¿qué sumas son
        alcanzables con bloques $(p, q)$, y cómo se compara este
        menú numerable con la carta completa de Riemann
        (pregunta 14)?
\end{enumerate}

\medskip
\noindent\textbf{Parte VI --- Epílogo: $\gamma$ en acción, y
síntesis.}
\begin{enumerate}[resume]
  \item Identifíquese la suma de la \omterm{def:b1:series:def}{serie} convergente
        $\sum_{k\geq1} \bigl(\frac1k - \ln\frac{k+1}{k}\bigr)$
        (\cref{exo:b1:series:7}): véase que vale $\gamma$.
  \item Ejecútese la receta de Riemann (pregunta 14) para el
        objetivo $t = 1$ y enumérense los doce primeros términos
        producidos ($1, \frac13, -\frac12, \frac15, -\frac14,
        \frac17, \frac19, -\frac16, \frac1{11}, \frac1{13},
        -\frac18, \frac1{15}$), calculando la suma parcial
        ($\approx 0.980$) --- obsérvese respirar al algoritmo en
        torno a su objetivo.
  \item Véase que alguna reordenación de la \omterm{def:b1:series:def}{serie} armónica alternada
        diverge a $+\infty$ \emph{(bloques de términos positivos lo
        bastante largos para ganar $1$ cada vez, usando la
        pregunta 12, separados por términos negativos sueltos)}.
  \item Afínese el~\cref{ex:b1:series:harmonicstack} con la ley de
        $\gamma$: véase que el primer índice con $H_N \geq 20$
        cumple $N = \eu^{\,20 - \gamma}\,(1 + o(1)) \approx
        2.7\cdot10^{8}$ --- la \omterm{pb:b1:series:1}{constante de Euler} es exactamente la
        corrección que le faltaba al encuadre tosco.
  \item Síntesis, una frase para cada punto: (i) la ley de $\gamma$
        y qué aporta cada una de sus tres piezas ($\ln n$, $\gamma$,
        $\frac{1}{2n}$); (ii) por qué la convergencia condicional
        hace que la suma dependa del orden mientras que la
        \omterm{thm:b1:series:absolute}{convergencia absoluta} lo prohíbe; (iii) cómo la fórmula
        $(p,q)$ fue un \emph{cálculo} con la ley de $\gamma$ y no
        una afirmación abstracta de existencia; (iv) dónde
        continúan estos hilos --- las \omterm{def:b1:series:def}{series} de potencias y los
        productos de \omterm{def:b1:series:def}{series} en el volumen del segundo año, y el
        problema del fin de semana del volumen del tercer año sobre
        la fórmula de Stirling, donde la misma contabilidad de suma
        frente a \omterm{thm:b1:integration:def}{integral} funciona a plena potencia.
\end{enumerate}
\end{problem}
